\newcommand{\speechpart}{Part A, slides 8--9}
\documentclass[11pt,a4paper]{article}
\usepackage[utf8]{inputenc}
\usepackage[T1]{fontenc}
\usepackage{lmodern}
\usepackage[margin=2.2cm]{geometry}
\usepackage{xcolor}
\usepackage{titlesec}
\usepackage{fancyhdr}
\usepackage[most]{tcolorbox}
\usepackage{microtype}
\usepackage{needspace}
\usepackage[hidelinks]{hyperref}
\definecolor{unicablue}{RGB}{45,58,104}
\definecolor{unicared}{RGB}{153,38,45}
\definecolor{tint}{RGB}{228,231,240}
\setlength{\parindent}{0pt}
\setlength{\parskip}{6pt}
\linespread{1.15}
\titleformat{\section}{\Large\bfseries\color{unicablue}}{}{0pt}{}
\titlespacing*{\section}{0pt}{12pt}{2pt}
\newcommand{\slide}[2]{\Needspace{10\baselineskip}\section{#1}\textcolor{unicared}{\small\bfseries #2}\par}
\newcommand{\cue}[1]{{\small\itshape\color{unicared}[#1]}}
\newcommand{\pause}{\textcolor{unicared}{\,/\!/\,}}
\newtcolorbox{say}{breakable,colback=white,colframe=unicablue,boxrule=0pt,leftrule=3pt,arc=0pt,
  left=10pt,right=6pt,top=4pt,bottom=4pt,fontupper=\large,before upper={\setlength{\parskip}{8pt}}}
\pagestyle{fancy}\fancyhf{}
\fancyhead[L]{\small\color{unicablue}Measuring Agglomeration --- \speechpart}
\fancyhead[R]{\small\color{unicablue}4 November 2026}
\fancyfoot[C]{\small\thepage}

\begin{document}
{\color{unicablue}\Huge\bfseries Speech: slides 8--9\par}
\vspace{2pt}
{\Large Part A, points (d)--(e)\par}
\vspace{4pt}
{\small About 2 minutes 15 seconds. \cue{red brackets} = what to point at or do; \pause\ = a short pause. If Part A is drawn and you present it alone, open with slides 1--7 (separate PDF).}

\slide{Slide 8 --- A spillover of 1/(sigma minus 1) is Helpman in disguise}{about 1 min 15 s}
\begin{say}
\cue{title} Point (d): Helpman 1998, in steps. In Helpman each location produces many varieties. Free entry makes their number proportional to population, n-i equals kappa L-i, and each variety sells at a constant markup over marginal cost.

\cue{formula in the white card} Step one: the contribution of location i to the price index of n is the number of varieties times their price to the one minus sigma. Since the number of varieties is proportional to L, that is \textbf{L-i to the power one}, times a term without L. Step two: in Armington with the spillover the price is w over A-bar L to the lambda, so the contribution is \textbf{L-i to the lambda times sigma minus one}, times a term without L.

\cue{formula lambda = 1/(sigma-1)} Step three: \pause\ the two coincide when lambda times sigma minus one equals one. So \textbf{lambda equals one over sigma minus one}. \cue{the two cards} With sigma equal to 5, the value in ARM1, that is 0.25 --- the value we used in points (b) and (c). With sigma equal to 4, the value in ARM2 and in Redding and Sturm, it is \textbf{one third}: keep this number in mind for Part B.

\cue{bottom box, look at the Professor} Step four: an estimated lambda need not be a technological externality, like knowledge spillovers. It can be the \textbf{variety gains} of monopolistic competition: more people, more firms, more varieties, a lower price index. Same equations, different stories --- so the interpretation of lambda depends on the model behind it.
\end{say}

\slide{Slide 9 --- What would identify lambda?}{about 1 min}
\begin{say}
\cue{title} Point (e). If one cross-section cannot identify lambda, because the fundamentals adjust to rationalise any value, we need \textbf{variation over time}: a panel.

\cue{red card, ``What must change''} What must change? \textbf{Population}, or the density of employment around a location, for reasons unrelated to productivity: a border that closes, a country divided, a new infrastructure.

\cue{blue card, ``What must stay fixed''} What must stay fixed? The \textbf{fundamentals}: the shock must move people without changing what each place intrinsically is.

\cue{formula} Then we invert the model at both dates. \pause\ With A-bar fixed, the change in log A is lambda times the change in log L, and that relation identifies lambda.

\cue{bottom box} This is exactly what the next two parts do: the division of Germany shocks market access across cities, and the Berlin Wall takes away the employment around the blocks near the old centre without changing what those blocks are.

\cue{hand over} My colleague will now present Part B.
\end{say}

\vspace{6pt}
\begin{tcolorbox}[colback=tint,colframe=tint,arc=0pt]
\small\textbf{Numbers to remember.} Helpman: exponent 1; Armington: $\lambda(\sigma-1)$ \ $\cdot$\ $\lambda = 1/(\sigma-1)$: 0.25 at $\sigma=5$ (ARM1), 1/3 at $\sigma=4$ (ARM2, Redding--Sturm) \ $\cdot$\ $\Delta\log A = \lambda\,\Delta\log L$ with $\bar A$ fixed.

\textbf{Words to get right.} \emph{one over sigma minus one}, never ``one over one minus sigma'' \ $\cdot$\ $\sigma=5$ vs 4 is ARM1 vs ARM2, not Armington vs Helpman \ $\cdot$\ the shock moves \emph{population}; productivity moves as a consequence.
\end{tcolorbox}
\end{document}
